\documentclass[10pt]{article}
\usepackage[T1]{fontenc}
\usepackage[utf8]{inputenc}
\usepackage[french]{babel}
\usepackage{amsmath,amssymb,mathrsfs,mathpazo}
\usepackage{geometry}
\geometry{a4paper,margin=2cm}
\usepackage{graphicx,array,tabularx,tabularray,colortbl,multirow,multicol,booktabs,enumitem,pifont,ulem,fancybox,diagbox,cancel}
\usepackage{pstricks,pstricks-add,pst-plot,pst-tree,pst-node,pst-text,pst-func,pst-math,pst-bspline,pst-3dplot,pst-eucl,pst-solides3d}
\let\pscalculate\undefined
\usepackage{pst-bezier}
\usepackage{tikz,pgfplots,tkz-tab}
\usetikzlibrary{arrows,arrows.meta,patterns,calc,positioning,trees,decorations.markings,shapes.geometric,shapes.misc,shapes.arrows,intersections,decorations.pathmorphing,decorations.pathreplacing,backgrounds,shadows,babel}
\pgfplotsset{compat=1.18}
\usepgfplotslibrary{fillbetween}
\usepackage[np]{numprint}
\usepackage{listings,makecell,pas-tableur,twoopt}
\let\barre\undefined
\usepackage[pstricks]{tablvar}
\definecolor{darkgreen}{RGB}{0,100,0}
\providecommand{\red}{\color{red}}
\providecommand{\blue}{\color{blue}}
\providecommand{\green}{\color{green}}
\providecommand{\black}{\color{black}}
\providecommand{\white}{\color{white}}
\providecommand{\gray}{\color{gray}}
\providecommand{\cyan}{\color{cyan}}
\providecommand{\magenta}{\color{magenta}}
\providecommand{\yellow}{\color{yellow}}
\providecommand{\violet}{\color{violet}}
\providecommand{\orange}{\color{orange}}
\providecommand{\euro}{\texteuro}
\providecommand{\EUR}{\texteuro}
\providecommand{\ieme}{\textsuperscript{e}}
\providecommand{\degres}{\ensuremath{^\circ}}
\providecommand{\degree}{\ensuremath{^\circ}}
\providecommand{\vv}[1]{\overrightarrow{#1}}
\providecommand{\Coord}[1]{#1}
\providecommand{\decosix}{}
\providecommand{\starredbullet}{\ensuremath{\bullet}}
\pagestyle{empty}
\setlength{\parindent}{0pt}
\usepackage{adjustbox}\DeclareUnicodeCharacter{2236}{:}\DeclareUnicodeCharacter{2190}{\ensuremath{\leftarrow}}
\providecommand{\R}{\mathbb{R}}
\providecommand{\N}{\mathbb{N}}
\providecommand{\Z}{\mathbb{Z}}
\providecommand{\Q}{\mathbb{Q}}
\providecommand{\D}{\mathbb{D}}
\providecommand{\C}{\mathbb{C}}
\providecommand{\e}{{\mathrm{e}}}
\providecommand{\vect}[1]{\overrightarrow{#1}}
\providecommand{\vectt}[1]{\overrightarrow{\mathrm{#1}}}
\providecommand{\vvt}[1]{\overrightarrow{\mathrm{#1}}}
\providecommand{\np}[2][]{{#2~#1}}
\providecommand{\vv}[1]{{\overrightarrow{#1}}}
\renewcommand{\d}{{\mathrm{d}}}
\providecommand{\degres}{{^\circ}}
\providecommand{\degre}{{^\circ}}
\providecommand{\blue}{\color{blue}}
\providecommand{\red}{\color{red}}
\providecommand{\enskip}{\quad}
\providecommand{\llbracket}{\lbrack\!\lbrack}
\providecommand{\rrbracket}{\rbrack\!\rbrack}
\providecommand{\euro}{{\text{€}}}
\providecommand{\no}{{\text{n°}}}
\providecommand{\interval}[2]{{[#1;#2]}}
\providecommand{\intervFO}[2]{{[#1;#2[}}
\providecommand{\dotfill}{\ldots}

\providecommand{\ds}{\displaystyle}
\providecommand{\textding}[1]{#1}
\providecommand{\Oij}{\(\mathrm{O};\vec{\imath},\vec{\jmath}\)}
\providecommand{\Oijk}{\(\mathrm{O};\vec{\imath},\vec{\jmath},\vec{k}\)}
\providecommand{\Ouv}{\(\mathrm{O};\vec{u},\vec{v}\)}
\providecommand{\intervFF}[2]{\ensuremath{[#1;#2]}}
\renewcommand{\euro}{\ifmmode\text{\texteuro}\else\texteuro\fi}
\DeclareUnicodeCharacter{00B0}{\ifmmode^{\circ}\else\textdegree\fi}
\DeclareUnicodeCharacter{20AC}{\ifmmode\text{\texteuro}\else\texteuro\fi}
% Énoncés : Éducation nationale. Transcription et corrigés : APMEP, Denis Vergès et contributeurs. Découpage et classement : CoopMaths, Sébastien Lozano, Éric Elter et contributeurs. Adaptation de la présentation : Maths Explorer.
% CC BY-SA : https://coopmaths.fr/www/about/
% https://coopmaths.fr/alea/static/e3c/2020/tex/e3c_2020_00_sujet16_3.tex
\begin{document}
\noindent\textit{APMEP / CoopMaths. CC BY-SA. Présentation adaptée par Maths Explorer.}\par\medskip
\medskip

Une chaîne de salons de coiffure propose à ses \np{5000} clients qui viennent pour une
coupe deux prestations supplémentaires cumulables :

\begin{itemize}
\item  une coloration naturelle à base de plantes appelée \og couleur-soin \fg,
\item des mèches blondes pour donner du relief à la chevelure, appelées \og effet
coup de soleil \fg.
\end{itemize}

Il apparaît que \np{2000} clients demandent une \og couleur-soin \fg. Parmi ceux qui ne
veulent pas de \og couleur soin \fg, $900$ demandent un \og effet coup de soleil \fg. Par
ailleurs, $650$ clients demandent une \og couleur soin \fg et un \og effet coup de soleil \fg.

On notera $C$ l'évènement \og le client souhaite une \og couleur-soin \fg.

On notera $E$ l'évènement \og le client souhaite un \og effet coup de soleil \fg.

\begin{enumerate}
\item Recopier sur votre copie et compléter le tableau suivant :

\begin{center}
\begin{tabularx}{0.7\linewidth}{|*{3}{>{\centering \arraybackslash}X|}m{0.75cm}|}
\hline
				& $C$ 	&$\overline{C}$	& Total\\\hline
$E$				&		& $900$			&\\\hline
$\overline{E}$ 	&		&				&\\\hline
Total			&		&				&\np{5000}\\\hline
\end{tabularx}
\end{center}
\smallskip

\item On interroge un client au hasard parmi les \np{5000} clients.
	\begin{enumerate}
		\item Quelle est la probabilité qu'il ait choisi les deux prestations : \og couleur soin \fg{}
et \og effet coup de soleil \fg ?
		\item Calculer $P_E\left(\overline{C}\right)$.
	\end{enumerate}
\item On a des prix différents suivant la prestation fournie. On appelle $X$ le prix payé
en euros par chaque client.

\begin{tabularx}{\linewidth}{|m{1.85cm}|m{2cm}|*{3}{>{\centering \arraybackslash}X|}}
\hline
&Coupe seule&Coupe avec \og couleur soin \fg &Coupe avec \og effet coup desoleil \fg &\small Coupe avec\og couleur soin \fg{} et \og effet coup de soleil \fg \\\hline
Valeurs de $k$ en \euro 	&\centering  20	& 50 	&65 	&80\\\hline
$P(X = k)$					&		&		& 0,18	&0,13\\
\hline
\end{tabularx}

\smallskip

Après avoir recopié et complété le tableau, calculer l'espérance de $X$.
\end{enumerate}

\vspace{0,5cm}

\end{document}
